% Template: artigo-basico
% Um artigo científico standalone em ABNT-like, sem dependências.
%
% Derivado/inspirado em github.com/Eranot/Unotex (LPPL 1.3c).
% Simplificado para uso como standalone (sem unotex.cls).

\documentclass[12pt,a4paper,oneside]{article}
\usepackage[utf8]{inputenc}
\usepackage[T1]{fontenc}
\usepackage[brazilian]{babel}
\usepackage[left=3cm,right=2cm,top=3cm,bottom=2cm]{geometry}
\usepackage{indentfirst}
\usepackage{setspace}
\usepackage{titlesec}

% ABNT-like: espaçamento 1.5, parágrafo com recuo, sem paragraph skip.
\onehalfspacing
\setlength{\parindent}{1.25cm}
\setlength{\parskip}{0pt}

% Título de seções em CAIXA ALTA.
\titleformat{\section}{\normalfont\large\bfseries\MakeUppercase}{\thesection}{1em}{}

\title{\textbf{ {{ titulo | default(value="Título do Artigo") }} }}
\author{ {{ autor | default(value="Nome do Autor") }} \\
         \small\texttt{ {{ email | default(value="autor@exemplo.br") }} } }
\date{ {{ data | default(value="\today") }} }

\begin{document}

\maketitle

\begin{abstract}
    {{ resumo | default(value="Resumo do trabalho aqui: contexto, objetivo, método, resultados e conclusão em até 250 palavras.") }}

    \vspace{6pt}\noindent\textbf{Palavras-chave:} {{ palavras_chave | default(value="palavra-1; palavra-2; palavra-3") }}.
\end{abstract}

\section{Introdução}

{{ introducao | default(value="Contextualize o problema, apresente a motivação e enuncie o objetivo do trabalho.") }}

\section{Desenvolvimento}

{{ desenvolvimento | default(value="Descreva método, dados, análise. Divida em subseções conforme necessário.") }}

\section{Considerações Finais}

{{ conclusao | default(value="Retome o objetivo, sintetize os principais achados e aponte trabalhos futuros.") }}

\begin{thebibliography}{9}
    \bibitem{eranot2020} ERANOT. \textit{Unotex: modelo LaTeX para trabalhos acadêmicos da Unochapecó}. GitHub, 2020. Disponível em: \url{https://github.com/Eranot/Unotex}. Acesso em: \today.
\end{thebibliography}

\end{document}
